\begin{table}[tbp]
\caption{Research questions and observation boundaries in the Redis study.}
\label{tab:ch9-design}
\centering\small\renewcommand{\arraystretch}{1.14}
\begin{tabularx}{\linewidth}{@{}>{\raggedright\arraybackslash}p{12mm}>{\raggedright\arraybackslash}X>{\raggedright\arraybackslash}p{45mm}@{}}\toprule
RQ & Analytical object & Observation window \\\midrule
RQ1 & Redis repository metrics. & October 2023--January 2025.\\
RQ2 & Counts and OpenRank of operational contributor groups. & October 2023--January 2025.\\
RQ3 & Annual top-ten contributors and attributed contribution shares. & 2017--June 2025; final year incomplete.\\
RQ4 & Redis/Valkey participation and contributor connections. & Monthly comparison: October 2023--January 2025.\\
RQ5 & Unique email domains as a proxy for organizational variety. & October 2023--January 2025.\\
RQ6 & Elasticsearch and Terraform comparison series. & July 2020--July 2021; February 2023--February 2024.\\\bottomrule
\end{tabularx}
\tablenote{Adapted from \citet[Sections 4--5]{han2026redis}. Question labels and periods refer to the source analyses. The annual selection, monthly groups, and two-project series are not one common longitudinal cohort.}
\end{table}
